\begin{center}\begin{tikzpicture}[font=\small,>=stealth,line join=round]
\draw[gray!40,rounded corners=3pt] (0,0) rectangle (14.2,9.0);
\node at (7.1,8.4) {例4：输入能成对，不等于输出符合奇偶规律};
\begin{scope}[shift={(3.6,2.9)},x=1.15cm,y=1.0cm]\draw[->,gray] (-2.2,0)--(2.4000000000000004,0) node[right]{$x$};\draw[->,gray] (0,-0.6)--(0,4.7) node[left]{$y$};\begin{scope}\clip (-2.2,-0.6) rectangle (2.2,4.5);\draw[blue,thick,domain=-2.2:2.2,samples=120] plot(\x,{\x*\x+\x});\end{scope}\fill[orange!85!black] (1,2) circle(2pt) node[right]{$(1,2)$};\fill[orange!85!black] (-1,0) circle(2pt) node[below left]{$(-1,0)$};\draw[orange!85!black,dotted] (-1,0)--(-1,2)--(1,2);\end{scope}\node[align=left,anchor=west] at (7.6,4.8) {$f(x)=x^2+x$\\[12pt]$D=\mathbb R$ 对称\\[12pt]$f(1)=2,\ f(-1)=0$\\[12pt]两值既不相等，也不互为相反数\\[12pt]$f(0)=0$ 也不能说明它是奇函数};
\node[text width=12.8cm,align=center] at (7.1,.65) {找到一对违例点，就能否定全称结论。};
\end{tikzpicture}\end{center}
